\begin{table}[H]
\centering
\caption{Simulated Method of Moments: Target Moments and Model Fit}
\begin{threeparttable}
\begin{tabular}{lccc}
\toprule
Moment & Data & Model & Difference \\
\midrule
GR employment response ($h=12$) & $-0.0435$ & $-0.0526$ & $-0.0091$ \\
GR employment response ($h=48$) & $-0.0527$ & $-0.0424$ & $0.0103$ \\
GR employment response ($h=84$) & $-0.0507$ & $-0.0309$ & $0.0199$ \\
COVID recovery month & $18.0000$ & $10.0000$ & $-8.0000$ \\
Steady-state unemployment rate & $0.0550$ & $0.0588$ & $0.0038$ \\
Steady-state job-finding rate & $0.4000$ & $0.3681$ & $-0.0319$ \\
\midrule
$J$-statistic & \multicolumn{3}{c}{$22.20$ ($df=2$, $p=0.0000$)} \\
Number of moments & \multicolumn{3}{c}{6} \\
Number of parameters & \multicolumn{3}{c}{4} \\
\bottomrule
\end{tabular}
\begin{tablenotes}[flushleft]
\small
\item \textit{Notes:} The table reports the fit of the Simulated Method of Moments (SMM) estimation. Four parameters ($\lambda$, $\chi_1$, $A$, $\kappa$) are estimated by minimizing the weighted distance between six target moments and their model counterparts; $\rho_a = 0.99$ is externally calibrated. Model employment responses are scaled to data units using the peak employment ratio (data peak / model peak). The $J$-statistic tests the overidentifying restrictions (6 moments $-$ 4 parameters $=$ 2 degrees of freedom). Standard errors are computed via parametric bootstrap (200 draws) with perturbations drawn from the estimated LP standard errors.
\end{tablenotes}
\end{threeparttable}
\label{tab:smm}
\end{table}
